% C.3bk finalized after complete root and both peer proof/source audits.
\paragraph{C.3bk (Proved): integral higher kernels from products of cubic surfaces.}
Let $k$ be an algebraically closed field of characteristic $p>0$.
For each of the following pairs $(r,N)$ there exist an excellent normal
strictly henselian local $k$-domain $R$ of dimension $N$, with residue
field $k$, and $0\ne f\in\mathfrak m_R$, such that the puncture $X$
and its principal punctured divisor $Y=X\cap V(f)$ are regular and
\[
 \ker\bigl(\mathrm{CH}^r(X)_{\rm alg}
             \longrightarrow\mathrm{CH}^r(Y)_{\rm alg}\bigr)
 \text{ contains a class of exact order }p:
 \qquad
 \begin{cases}
 r=2,&N\ge4,\\
 r=3,&N\ge7,\\
 r\ge4,&N\ge3r-1.
 \end{cases}
 \tag{C.3bk.1}
\]
All algebraic equivalences are realized on actual fixed products with
smooth proper parameter curves. The data can be chosen over
$\overline{\mathbb F}_p$ and then extended to $k$. In particular the
new $r=2,N=4$ assertion holds even over $\overline{\mathbb F}_p$;
C.3ay in that dimension used a specified transcendental coefficient
field. These are integral order-$p$ classes, and give no nonzero
rational kernel. The proof makes no assertion for $r\ge4$ in the
remaining high interval $2r+1\le N<3r-1$.

\begin{proof}
\textbf{1. Specified projective factors and their coherent cohomology.}
Put $s=r-1$ and $m=N-2r$. In the three ranges of (C.3bk.1) one has
$m\ge0$, with $m=0$ used only when $r=2$.
Choose the ordinary plane quartic $C=(F=0)$, quartic $H$, and
nondegenerate normal-deformation form $Q_H$ from C.3bg and C.3bj.
Their finite-field enumeration gives actual data in every
characteristic. Let $W_1,\ldots,W_s$ be copies of the specified
smooth cubic surface of C.3bj, with hyperplane line bundles $A_i$
and lines $D_i$. In characteristic different from three these can
be the Fermat cubic with the stated line; in characteristic three
use $u s^2+v t^2+u^2t+v^2s=0$ and $D_i=(u=v=0)$.
Write
\[
 \begin{gathered}
 T_0=\prod_{i=1}^s W_i,\qquad S=T_0\times C,\qquad
 B=\Bigl(\boxtimes_{i=1}^s A_i\Bigr)\boxtimes K_C,\\
 V=\mathbb P^m\times S,\qquad
 L=\mathcal O_{\mathbb P^m}(1)\boxtimes B,
 \qquad \dim V=m+2r-1=N-1.
 \end{gathered}
 \tag{C.3bk.2}
\]
When $m=0$, the projective-space factor and its line bundle are
understood to be a point and its trivial line bundle.

For every $i$, the characteristic-independent calculation in
C.3bj gives the rank-two map
\[
 \mu_i:H^0(W_i,A_i)\longrightarrow
     H^1(W_i,A_i\otimes\Omega^1_{W_i/k}),
 \qquad t\longmapsto t\,d\log\mathcal O_{W_i}(D_i).
 \tag{C.3bk.3}
\]
Its dual form on the line is the functional on binary quadratics
annihilating the two normal quadratics $Q_0,Q_1$. They have no common
zero, so the induced two-dimensional bilinear form is nondegenerate,
also in characteristic two. This is the coefficient-one normal-sequence
calculation of (C.3bj.4)--(C.3bj.5); no factorial is introduced when
these maps are tensored.

The cubic hypersurface sequence in $\mathbb P^3$ gives
$H^j(W_i,A_i^n)=0$ for $j>0,n\ge0$. For example the only possible
top group is dual to $H^0(W_i,A_i^{-n-1})$, which is zero.
On the curve, $H^1(K_C)=k$ and $H^1(K_C^n)=0$ for $n\ge2$.
K\"unneth and projective-space cohomology consequently give
\[
 \begin{aligned}
 U_0:=H^1(V,L)
   &=H^0(\mathbb P^m,\mathcal O(1))
       \otimes\bigotimes_{i=1}^s H^0(W_i,A_i)
       \otimes H^1(C,K_C),\\
 H^j(V,L^n)&=0&& (j\ge2,\ n\ge0),\\
 H^1(V,L^n)&=0&& (n\ge2).
 \end{aligned}
 \tag{C.3bk.4}
\]
The same calculation for $n=0$ gives
$\operatorname{Pic}^0(V)=J:=\operatorname{Pic}^0(C)$ and
$H^1(V,\mathcal O_V)=H^1(C,\mathcal O_C)$.

\textbf{2. A smooth incidence section with a nonzero product cup.}
For $m\ge1$, choose sections $a_0,\ldots,a_m$ of the very ample
line bundle $B$ so that their common zero scheme is smooth of the
expected codimension when $m+1\le\dim S$, and is empty when
$m+1>\dim S$. Successive Bertini sections give a nonempty open set
of such tuples; after a transverse zero-dimensional intersection,
the next section can be chosen to avoid all its points. Additional
sections can be zero once the base locus is empty. The applicable
primary input is the very-ample immersion version of Bertini
\citep[Lemma 33.47.3, Tag 0FD6]{D_StacksBertiniUniformHigher}.
Put
\[
 \ell=\sum_{i=0}^m x_i a_i=\sum_{j=0}^2 g_jz_j,
 \qquad
 g_j\in H^0(\mathbb P^m,\mathcal O(1))
                    \otimes\bigotimes_iH^0(W_i,A_i).
 \tag{C.3bk.5}
\]
We also require the three $g_j$ to be independent and their span to
have nonzero image under $1\otimes\bigotimes_i\mu_i$.
These rank conditions define a nonempty open set in the same affine
parameter space. To see nonemptiness, choose $s_i,t_i$ whose images
under $\mu_i$ are independent. The three tensors
\[
 x_0\prod_i s_i,\qquad
 x_0t_1\prod_{i>1}s_i,\qquad
 x_1\prod_i s_i
\]
are independent and have nonzero product cup. Irreducibility of the
parameter space makes this open intersect the Bertini open.
For the sole case $m=0,r=2$, choose a general section of $A_1\boxtimes K_C$.
Its three coefficients $g_j\in H^0(W_1,A_1)$ can be independent,
and their span cannot lie in the two-dimensional kernel of $\mu_1$.
These conditions again meet the nonempty smooth-divisor open.

In every case the divisor $T=(\ell=0)\subset V$ is smooth.
For $m\ge1$, away from the common base locus a projective-space
derivative is nonzero. At a base point, the differentials of the
sections are independent, so every nonzero linear combination with
coefficients $(x_i)$ is nonzero. For $m=0$ this is the Bertini
smoothness already imposed. All choices lie in nonempty opens defined
over finite fields. Enumerating finite extensions and coefficient
tuples, and selecting the first satisfying the Jacobian and rank
tests, specifies them by a terminating procedure. No bound on its
running time is being asserted.

\textbf{3. The actual strict-local deformation and its split cut.}
In the total space of
\[
 \mathcal O_{\mathbb P^m}(-1)
       \boxtimes\Bigl(\boxtimes_i A_i^{-1}\Bigr)
       \boxtimes\mathcal O_{\mathbb P^2}(-1)
\]
over $\mathbb P^m\times T_0\times\mathbb P^2$, take
\[
 Z_0:\quad F+\tau\ell H=0.
 \tag{C.3bk.6}
\]
The base factors are embedded by their full degree-one section rings,
so the product Segre contraction is projective and is an isomorphism
away from its zero section. Contract $Z_0$, take the reduced affine
image and its normalization, localize at its vertex, and strictly
henselize; call the resulting local ring $R$. The same graph proof
as in C.3bj applies to any number of the smooth cubic factors.
Indeed $F$ does not divide $\ell H$, so the equation, linear in
$\tau$, is integral. Its differential is nonzero along the exceptional
$C$. The exceptional fibre $V$ is connected, so the finite normalization
has only one point over the vertex. The image of the nonregular locus
outside a regular neighbourhood of $V$ is closed by properness and
avoids that vertex; it disappears after localization. Thus
\[
 \mathcal Z\longrightarrow\operatorname{Spec}R
\]
is projective, has regular source, and is an isomorphism away from
its exceptional divisor $V$. The Segre generators give
\[
 \dim R=N,\qquad \mathfrak m_R\mathcal O_{\mathcal Z}=I_V,
 \qquad I_V/I_V^2=L.
 \tag{C.3bk.7}
\]
This also fixes the actual first neighbourhood $E:=V_2$.
Every projection to $W_i$ is retained, so
$\mathcal D_i:=\mathcal O_{W_i}(D_i)$ is an actual line bundle
on $\mathcal Z$. Take $f=\tau\ell$, the corresponding linear
combination of the Segre coordinates. On its strict transform the
equation becomes exactly $F=0$. That transform is the split total
space of $L|_T^{-1}$, with zero section $T$, followed by the same
strict-local base change. In particular the puncture $X$ and its
principal cut $Y$ are regular.

\textbf{4. An actual ordinary Picard endpoint.}
The first-neighbourhood Picard extension is
\[
 0\longrightarrow U_0\longrightarrow G\longrightarrow J\longrightarrow0.
 \tag{C.3bk.8}
\]
The transition-unit/Atiyah calculation of C.3bg identifies its
extension-class map with
$s\mapsto\sum_jQ_H(z_j,s)g_j$.
Its image is the three-dimensional space
$U=\ell H^1(V,\mathcal O_V)=\langle g_0,g_1,g_2\rangle$.
This uses precisely the characteristic-independent obstruction formula
\citep[Theorem 2.6 and formula (2.7.2), pp.~166--167]{D_Illusie1972Torsors}.
Ordinarity of $J$ gives its etale Verschiebung
$v:J^{(p)}\to J$, with $K=\ker v\simeq(\mathbb Z/p)^3$ and
$v^*H^1(J,\mathcal O_J)=0$.
The integral extension classification
\citep[Section 3.2.3, equation (3.8), p.~23]{D_BrionIsogenyII}
therefore gives a homomorphism $b:J^{(p)}\to G$ over $v$.
The Hom--Ext sequence identifies $b(K)$ with an etale subgroup
spanning $U$, and $b|_K$ is injective.

The product of the divisor cups is nonzero on $U$ by Step 2.
Its component is the direct coherent K\"unneth summand
\[
 H^0(\mathbb P^m,\mathcal O(1))\otimes
       \bigotimes_{i=1}^s H^1(W_i,A_i\otimes\Omega^1_{W_i/k})
       \otimes H^1(C,K_C)
 \subset H^r(V,L\otimes\Omega^{r-1}_{V/k}).
 \tag{C.3bk.9}
\]
Consequently one can choose $\eta\in K$ so that, for $a=b(\eta)$,
\[
 a\smile\prod_{i=1}^s d\log\mathcal D_i\ne0.
 \tag{C.3bk.10}
\]
The tensor product here has its fixed factor order; its coefficient
is one, up to sign, rather than $(r-1)!$.

The divisor sequence and (C.3bk.4) imply
\[
 H^1(T,L|_T)=U_0/U,\qquad
 H^1(T,L^n|_T)=H^2(T,L^n|_T)=0\quad(n\ge2).
 \tag{C.3bk.11}
\]
The first target neighbourhood is split. Its Picard functor splits
into that of $T$ and the vector group $U_0/U$; no identification
of the whole $\operatorname{Pic}^0(T)$ with $J$ is required.
The restriction of $b$ to the first factor factors through $v$.
Its map to the vector factor is zero because it is a morphism from
an abelian variety, normalized to zero at the identity. Thus $a$
has zero first jet on the target, and (C.3bk.11) makes all its
higher endpoint lifts unique.

For clarity we describe the actual family, rather than infer it
from the representability of a local Picard functor. Choose a smooth
proper connected curve $P\subset J^{(p)}$ through $0$ and $\eta$.
Pull the normalized Poincare bundle to $V\times P$ using $v|_P$.
Its first mixed Atiyah obstruction is killed by
$v^*:H^1(J,\mathcal O_J)\to H^1(J^{(p)},\mathcal O)$,
so it lifts to $E\times P$. Its map to $G$ differs from $b|_P$
by a constant map to $U_0$, removed by tensoring with its inverse
fibre at zero. For $n\ge2$, K\"unneth, (C.3bk.4), and
$H^2(P,\mathcal O_P)=0$ give $H^2(V\times P,L^n)=0$.
Thus this line bundle lifts compatibly through every exceptional
thickening. Proper formal existence gives a line bundle on
$\mathcal Z_{\widehat R}\times P$.
Finite-data Popescu descent, including the inverse line bundle and
the first-jet identification, followed by smooth henselian lifting
modulo $\mathfrak m_R^2$, gives an actual line bundle on
$\mathcal Z\times P$ with that first jet, exactly as in C.3bf--C.3bj.
Normalize the fibre at zero to be trivial and call the fibre at
$\eta$ by $\mathcal A$.
The uniqueness in (C.3bk.4) and (C.3bk.11), proper formal existence,
and Elkik's punctured Picard comparison
\citep[unnumbered corollary, p.~573]{D_Elkik1973} prove
\[
 \begin{gathered}
 \mathcal A|_V=\mathcal O_V,\qquad
 \mathcal A|_E\text{ has relative class }a,\\
 \mathcal A|_X\in\operatorname{Pic}(X)_{\rm alg},\qquad
 (\mathcal A|_X)^p=\mathcal O_X,\qquad
 \mathcal A|_Y=\mathcal O_Y.
 \end{gathered}
 \tag{C.3bk.12}
\]

\textbf{5. Divisibility and the integral relative Milnor detector.}
For a smooth variety $Q$ write
$G^{a,b}(Q)=H^a(Q,\mathcal K_b^M)$.
The following integral divisibility holds for our $V$:
\[
 H^{r-1}(V,\mathcal K_r^M)\text{ is divisible}.
 \tag{C.3bk.13}
\]
Here are the complete generation steps, extending the degree-two
argument of C.3bj. Rost's localization and affine homotopy, with
pullback along the projection, give
\[
 G^{a,b}(Q)\oplus
       G^{a-1,b-1}(\mathbb P^{d-1}\times Q)
       \twoheadrightarrow G^{a,b}(\mathbb P^d\times Q).
 \tag{C.3bk.14}
\]
The first summand lifts restriction to the affine open, and the
second is the image supported on the hyperplane. For the blow-up
$\rho$ of a point-product centre, the difference
$c-\rho^*\rho_*c$ is zero off the exceptional divisor by proper/open
base change and the isomorphism there. It is therefore in the image
of $G^{a-1,b-1}(E)$, where $E=\mathbb P^1\times Q$.
The exact primary inputs are Rost's Proposition 4.1(3), localization
in Section 5, Corollary 6.5, Proposition 8.6, and the smooth pullbacks
in Section 12 \citep[pp.~351, 355--360, 370, 382--383]{cyRost1996}.
The integral Gersten comparison with the ordinary Milnor symbol
sheaf is \citet[Theorems 7.1 and 7.5]{cyKerz2009Milnor}; all residue
fields in use are infinite.
Each $W_i$ is the blow-up of $\mathbb P^2$ in six points
\citep[Chapter 4, Proposition 2.4, p.~178]{D_HuybrechtsCubics}.
Induction on the number of these factors, the projective-space
dimension, and the cohomological index in (C.3bk.14) reduces
$G^{a,a+1}(V)$ to quotients of direct sums of the curve groups.
For the curve these are $G^{0,1}(C)=k^*$,
\[
 G^{1,2}(C)=\operatorname{coker}
       \left(K_2^M(k(C))\longrightarrow\bigoplus_{c\in C(k)}k^*\right),
\]
and zero for $a\ge2$.
They are divisible; sums and quotients of divisible abelian groups
remain divisible. This proves (C.3bk.13).

On the common topological space of $E$ and $V$, put
\[
 \mathcal J_r=\ker(\mathcal K_r^M(E)\longrightarrow\mathcal K_r^M(V)).
\]
Units lift, and every downstairs Steinberg relation lifts by choosing
$u$ and $1-u$. Hence the Milnor map is onto, and its kernel is
locally generated by symbols having an entry $1+i$, $i\in I=L$.
Since $(1+i)^p=1$, multilinearity implies $p\mathcal J_r=0$ in
every characteristic. Local smoothness splittings give the canonical
short exact conormal sequence
\[
 0\longrightarrow I\xrightarrow d\Omega^1_{E/k}|_V
              \longrightarrow\Omega^1_{V/k}\longrightarrow0.
\]
The exact inputs are formal smoothness and the usual conormal sequence
\citep[Proposition 10.138.13 and Lemma 10.131.9]{frStacksFormalSmoothRelativeMilnor,D_StacksConormalUniformHigher};
left injectivity is checked on a local split algebra
$\mathcal O_V\oplus I$, rather than assumed for arbitrary
square-zero immersions. Rank one identifies the kernel on $r$th
exterior powers with $I\otimes\Omega^{r-1}_{V/k}$.
The ordinary $d\log$ of Milnor symbols therefore gives a canonical
map, independent of the chosen splitting,
\[
 \begin{aligned}
 \tau_r:\mathcal J_r&\longrightarrow I\otimes\Omega^{r-1}_{V/k},\\
 \{1+i,u_2,\ldots,u_r\}&\longmapsto
       i\,d\log\bar u_2\wedge\cdots\wedge d\log\bar u_r.
 \end{aligned}
 \tag{C.3bk.15}
\]
Symbols with more than one infinitesimal entry have zero image
because the conormal submodule has rank one. They are not being
declared zero as Milnor symbols in characteristic two.

The actual line bundles in (C.3bk.12) give a relative Cech symbol
class $\xi\in H^r(V,\mathcal J_r)$ whose image under $\tau_r$ is
exactly the nonzero class (C.3bk.10), up to the fixed sign of factor
order. This Cech class maps canonically to sheaf cohomology; no
acyclicity of a cover for $\mathcal J_r$ is assumed.
The long exact cohomology sequence has preceding map from
$H^{r-1}(V,\mathcal K_r^M)$. It is zero by (C.3bk.13), since its
target is annihilated by $p$. Consequently
\[
 H^r(V,\mathcal J_r)\hookrightarrow H^r(E,\mathcal K_r^M),
 \qquad \xi\ne0.
 \tag{C.3bk.16}
\]
Only the ordinary Milnor symbol sheaf on $E$ is used; no Bloch
formula or Milnor--Quillen comparison is asserted on that nonreduced
scheme.

\textbf{6. The entire exceptional image in the stated dimension ranges.}
We prove the required integral assertion
\[
 c_1(L):\mathrm{CH}^{r-1}(V)\longrightarrow\mathrm{CH}^r(V)
       \text{ is injective in the cases of (C.3bk.1).}
 \tag{C.3bk.17}
\]
For $m\ge r-1$, put $h=c_1(\mathcal O_{\mathbb P^m}(1))$ and
$c=c_1(B)$. The integral projective-bundle formula expresses
$\gamma=\sum h^i\gamma_i$ uniquely
\citep[Lemma 42.36.2, Tag 02TX]{cyStacksProject}.
If $m\ge r$, the successive coefficients of $h^r,h^{r-1},\ldots,h$
in $(h+c)\gamma=0$ force all $\gamma_i=0$.
If $m=r-1$, the top coefficient is
$b=\gamma_{r-1}\in\mathrm{CH}^0(S)=\mathbb Z$ and the recurrence gives
\[
 \gamma_i=(-c)^{r-1-i}b,\qquad b c^r=0.
 \tag{C.3bk.18}
\]
Multiplying by $c^{\dim S-r}$ and taking degree gives
$b\deg(c^{\dim S})=0$ in $\mathbb Z$. The line bundle $B$ is ample,
so this degree is a positive integer. Thus $b=0$ integrally,
including when the characteristic divides that integer.
This proves (C.3bk.17) for $N=m+2r\ge3r-1$.

We give the two extra endpoints without discarding torsion.
The integral point-blow-up and projective-bundle formulas, applied
after product with an arbitrary smooth scheme $Q$, give
\[
 \mathrm{CH}^*(W_1\times\cdots\times W_s\times Q)
   =\bigotimes_i\mathrm{CH}^*(W_i)\otimes_{\mathbb Z}\mathrm{CH}^*(Q).
 \tag{C.3bk.19}
\]
Indeed a point-product blow-up adds exactly the exceptional divisor
basis element tensored with $\mathrm{CH}(Q)$; its normal bundle
on the centre is the constant rank-two bundle. The self-intersection
and projection formulas determine all products. The precise integral
maps are \citet[Theorems 2.4 and 2.8(3)]{cyJiang2023}.
Thus (C.3bk.19) is proved by those explicit product blow-ups; it is
not a K\"unneth assertion for arbitrary projective varieties.
For each cubic, its basis $H,E_1,\ldots,E_6,z$ satisfies
\[
 H^2=z,\quad E_i^2=-z,\quad HE_i=E_iE_j=0\ (i\ne j),
 \qquad A=3H-\sum_iE_i,\quad A^2=3z.
 \tag{C.3bk.20}
\]

For $r=2,m=0$, write $\gamma=d+\lambda$ in
$\operatorname{Pic}(W)\oplus\operatorname{Pic}(C)$.
Equation $(A+K_C)\gamma=0$ gives
\[
 A\cdot d=0,\qquad d\boxtimes K_C+A\boxtimes\lambda=0.
\]
Taking degree in the curve factor yields
$4d+(\deg\lambda)A=0$ in the free cubic Picard lattice.
Pairing with $A$ gives $3\deg\lambda=0$ in $\mathbb Z$.
Hence $\deg\lambda=0$ and $d=0$. The class $A$ is primitive in
that lattice, because every exceptional coefficient is $-1$.
Thus $A\boxtimes\lambda=0$ implies $\lambda=0$ in the full
Picard group. This proves (C.3bk.17) also for $r=2,N=4$.

For $r=3,m=1$, put $a=A_1$, $b=A_2$ and $c=a+b+K_C$.
Write $\gamma=\gamma_0+h\gamma_1$. Then
$\gamma_0=-c\gamma_1$ and $c^2\gamma_1=0$.
By (C.3bk.19), uniquely write
$\gamma_1=d_1+d_2+\lambda$.
The two pure cubic summands of $c^2\gamma_1=0$ give
\[
 3d_2+2(a\cdot d_1)b=0,\qquad
 3d_1+2(b\cdot d_2)a=0.
 \tag{C.3bk.21}
\]
For the integers $x=a\cdot d_1$ and $y=b\cdot d_2$, pairing gives
$2x+y=x+2y=0$, whence $x=y=0$ and then $d_1=d_2=0$.
The remaining relation is
\[
 (3z_1\boxtimes1+2a\boxtimes b+3\,1\boxtimes z_2)
                         \boxtimes\lambda=0.
\]
The pure summands give $3\lambda=0$; the coefficient of any
$E_{1,i}\boxtimes E_{2,j}$ gives $2\lambda=0$.
Since $3-2=1$, this kills $\lambda$ including all Picard torsion.
Thus $\gamma=0$, proving (C.3bk.17) for $r=3,N=7$.
Together with the general range this covers exactly (C.3bk.1).

\textbf{7. A nonzero actual ordinary cycle in the restriction kernel.}
Define on the regular resolution
\[
 \beta=c_1(\mathcal A)\prod_{i=1}^s c_1(\mathcal D_i)
                  \in\mathrm{CH}^r(\mathcal Z).
 \tag{C.3bk.22}
\]
If $\beta|_X=0$, localization gives
$\beta=i_*\gamma$ for $\gamma\in\mathrm{CH}^{r-1}(V)$.
Its restriction to $V$ is zero since $\mathcal A|_V$ is trivial.
The self-intersection formula therefore gives
$0=i^*i_*\gamma=-c_1(L)\gamma$.
Step 6 implies $\gamma=0$, so $\beta=0$ on $\mathcal Z$.

This contradicts (C.3bk.16). To make the use of the integral
Bloch formula precise, descend the finitely many actual line bundles
and the finite hypothetical rational-equivalence relation $\beta=0$
to one pointed etale neighbourhood of the algebraic vertex.
Its base-changed resolution is regular on a neighbourhood of the
exceptional fibre; shrink the base using the proper image of the
nonregular locus. It is then a smooth finite-type scheme over the
perfect field $k$. The etale local map with unchanged residue field
preserves the quotient modulo the square of the maximal ideal, so
its exceptional first neighbourhood is exactly $E$ and the restricted
bundles are exactly those used in (C.3bk.15)--(C.3bk.16).
The integral Bloch formula of
\citet[Theorems 7.1 and 7.5]{cyKerz2009Milnor}, with its product
compatibility \citep[Section 3, pp.~557--558, diagram (22)]{cyBraunling2013},
sends the product of first Chern classes to the product of their
Milnor classes. Its vanishing would restrict to zero in
$H^r(E,\mathcal K_r^M)$, whereas that restriction is the nonzero
image of $\xi$. This contradiction proves $\beta|_X\ne0$.
No cohomology-limit theorem or Chow formula for $E$ is needed.

The actual family producing $\mathcal A$ in Step 4, multiplied by
the fixed divisor classes $\mathcal D_i$, proves
$\beta|_X\in\mathrm{CH}^r(X)_{\rm alg}$.
Equation (C.3bk.12) makes this class vanish on $Y$ and annihilates
it by $p$. Its nonzero value therefore has exact order $p$.
This proves (C.3bk.1).
\end{proof}

\noindent\emph{Consequences and limitations.}
Together with C.3aw and C.3ay, this proves that for codimension three,
when universality is taken over all algebraically closed coefficient
fields in the stated local scope, integral restriction is injective
exactly in the vacuous dimensions $N\le3$. The lower counterexamples
may require the specified transcendental field; the present
$N\ge7$ examples exist already over $\overline{\mathbb F}_p$.
For codimension two, C.3bj already gives the universal classification;
the new endpoint strengthens its fixed-field scope at $N=4$.
For $r\ge4$ the interval $2r+1\le N<3r-1$ is deliberately left
outside this proof. None of these order-$p$ examples changes a
rational injectivity statement.
